\begin{tabular}{p{3.2cm} p{5.4cm} p{3.0cm} p{4.4cm}}
\toprule
\textbf{Study} & \textbf{Short Description} & \textbf{N} & \textbf{Inter-rater Reliability} \\
\midrule
This Study & LLM MITI 4.2.1 scoring, 6 runs (1 benchmark + 5 stochastic reruns), gpt-5-nano; Change Talk \& Ambivalence arms only & 5466 session-dimension units (1369 sessions) & ICC 0.45 \\
\midrule
Moyers et al.\ (2016) & MITI 4.0 development study, 4 coders; four global ratings & 50 sessions & Mean global ICC 0.89 (range 0.87--0.91) \\
Paranjothy et al.\ (2017) & MITI 4.1, 2 raters; breastfeeding peer-support sessions & 5 double-coded sessions (16 total) & Pooled global-score $\kappa=0.49$ \\
Kramer Schmidt et al.\ (2019) & MITI 4.2.1 (Danish), 5 raters; four global ratings & 52 sessions & Mean global ICC 0.72 (range 0.50--0.84) \\
Kramer Schmidt et al.\ (2022) & MITI 4.2.1, four international rating teams; four global ratings & 52 / 12 / 13 / 20 sessions & Global ICC ranges by team: DK 0.49--0.84; Dresden 0.46--0.73; Munich 0.11--0.45; US 0.00--0.25 \\
Cohen et al.\ (2024) & MITI 4.2 MI-TAGS, 3 trainee annotators versus an expert; four global ratings & 23 calibration sessions & Mean annotator--expert Pearson $r=0.53$ (range 0.39--0.64) \\
Frey et al.\ (2025) & MITI 4.2.1, independent coding team; school-based coaching & 142 sessions; IRR subset not stated & Four global-item ICCs ranged 0.85--1.00 \\
Jha et al.\ (2026) & MITI 4.2, 2 raters; model and clinical transcripts & 100 transcripts & Mean global ICC 0.87 (range 0.86--0.90) \\
\bottomrule
\end{tabular}
\par\smallskip
\begin{minipage}{\linewidth}
\footnotesize\textit{Note.} The first row reports the reliability of our LLM-based scoring procedure, computed as ICC(2,1) (two-way random effects, absolute agreement, single-rater reliability) across six full-corpus scoring runs, pooled across the four MITI 4.2.1 global dimensions and restricted to the Change Talk and Ambivalence treatment arms. Literature rows are limited to MITI 4.x studies with a verified numeric inter-rater reliability estimate for the global ratings or a directly reported global summary. Statistics are not fully harmonized: studies use different ICC models and units; Paranjothy et al.\ report kappa; and Cohen et al.\ report Pearson correlations, which measure association rather than absolute agreement. Comparisons are therefore descriptive rather than like-for-like. For Moyers et al.\ (2016), Kramer Schmidt et al.\ (2019), Cohen et al.\ (2024), and Jha et al.\ (2026), the displayed mean is the unweighted arithmetic mean of the four published dimension-specific coefficients, not a re-estimated pooled coefficient. Frey et al.\ (2025) report only the range for the four global-item ICCs. The Danish team in Kramer Schmidt et al.\ (2022) overlaps the 2019 Danish sample; the 2022 row is included because it also reports three additional international rating teams. N is the inter-rater sample where stated; otherwise the total session count and the missing IRR-subset detail are shown.
\end{minipage}